\section{Results}

\subsection{Validation of Monte Carlo Sampler}
\label{sec:validationMC}
The most interesting direction is to treat a TSP tour as a statistical-mechanical microstate. 
Let $x$ be a microstate in the $(N-1)!/2$ dimensional state space.
The probability to visit a state $x$ is then given by the formula \ref{eq:prob_visit}
Now this quantity is generally not known, as the state space cannot be explored fully for even a small number of nodes.

Nevertheless, a system of 11 nodes can still be solved quite quickly, as it has only 1814400 possible unique permutations.
The 11 nodes were generated by using the random package from NumPy with the seed \textit{30121999} in the following form. 
\begin{equation}
    x/y \sim \mathcal{U}_{\mathbb{Z}}(\text{low}, \text{high} - 1) + \mathcal{N}\left(0, \frac{|\text{high}|}{2}\right)
    \label{eq:dist_cities}
\end{equation}
All numbers of unique permutations were computed, and the partition function \ref{eq:partition}, the expectation value \ref{eq:expect}, and the variance \ref{eq:variance} were evaluated in the temperature range (0.05;50) with a resolution of 0.01.
This node configuration will serve as an exactly solved system, which we can use to analyze the sampling of the algorithms presented in section \ref{sec:Methods}.  
The distribution of the nodes as well as the optimal distance $\text{min}(L(\tau))$ is shown in figure \ref{fig:Nodes_solved}. 
\begin{figure}[H]
    \centering
    \includegraphics[width = 0.7\textwidth]{cities1.pdf}
    \caption{Eleven Randomly distributed Nodes according to eq. \ref{eq:dist_cities} and optimal distance}
    \label{fig:Nodes_solved}
\end{figure}
As the system is solved fully, we can also calculate the probability of visiting the optimal state for each temperature with formula \ref{eq:prob_visit}.
In figure \ref{fig:prob_visit} the probabilities for visiting the minimal-, maximal-, and mean-length state $\tau$ were calculated.

\begin{figure}[H]
    \centering
    \includegraphics[width= 0.7\textwidth]{prob_visit_optimal.pdf}
    \caption{top: $p_L(\tau|T)$ for different lengths produced by the Node configuration, bottom: length distribution}
    \label{fig:prob_visit}
\end{figure}
It is apparent that visiting the optimal state is most probable at low temperatures as predicted by the equilibrium Boltzmann distribution.
As $T\rightarrow 0^+$ only the shortest tour survives. Factoring out the weight of the shortest state from eq. \ref{eq:prob_visit} gives:
\begin{equation}
    p_L(\tau_{\mathrm{min}}|T) = \frac{1}{1 + \sum_{\tau \neq \tau_{\mathrm{min}}} e^{-(L(\tau)-L( \tau_{\mathrm{min}}))/T}}
    \label{eq:lowtemplimit}
\end{equation}
From eq. \ref{eq:lowtemplimit} follows that every difference $L(\tau)-L( \tau_{\mathrm{min}})$ approaches zero; therefore, the probability as a whole approaches 1.
Now this is nice, but what happens in the high-temperature limit?
\begin{figure}[H]
    \centering
    \includegraphics[width = 0.7\textwidth]{probabilities_colormap2.pdf}
    \caption{probabilities of reaching optimal length}
    \label{fig:hightemplimit}
\end{figure}
In figure \ref{fig:hightemplimit} the probabilities for lengths in the range of the lengths distribution from \ref{fig:prob_visit} are plotted up to a temperature of 10000.
At high temperatures every state change becomes equally probable, this is easy shown when taking the limit $T \rightarrow \infty$ of eq. \ref{eq:prob_visit}:
\begin{equation}
    p_L(\tau_i | T)_{T  \rightarrow \infty} = \frac{e^0}{\sum^{N_{\tau}} e^0}=\frac{1}{N_{\tau}}
\end{equation}

\subsubsection{Metropolis-Hastings Algorithm (Annealing)}
To test if the standard annealing algorithm truly samples the correct Boltzmann distribution, the exact solution was compared to the observables measured in between sweeps at constant temperature.
A sweep in this case consists of $N_{\mathrm{Nodes}}^2$ state change proposals.
This ensures that the system is equilibrated before measuring.
The implemented algorithm performed 8000 sweeps for each temperature in the range 0.05 to 50 with a resolution of 0.01.
However, 20 \% of the sweeps for each temperature were used for the warmup procedure to ensure that all states can be reached.
This is especially important for low temperatures, because the probability for state change is extremely low.
The first 2000 sweeps were performed with the temperature function: \\
\begin{equation}
    T(n_{\mathrm{sweep},i},T_{\mathrm{finish}}) = \frac{1200}{n_{\mathrm{sweep},i} + 1} + T_{\mathrm{finish}}
\end{equation}
Following the warm-up procedure, the algorithm was run with constant temperature, and after each sweep, the last length as well as the acceptance rate were observed and tracked.
This gave a sample set of 6000 for each temperature. The quantities in figure \ref{fig:annealing_solved_system} are then calculated using the standard formulas for the mean and variance.
The exact solution was calculated using the equations [\ref{eq:partition},\ref{eq:expect},\ref{eq:variance}, \ref{eq:heat_cap}].
\begin{figure}[H]
    \centering
    \includegraphics[width = 0.8\textwidth]{solved_system_0.5_50_0.01_plot.pdf}
    \caption{Comparison between the exact solution (red) and the sampled observables using the annealing algorithm (blue)}
    \label{fig:annealing_solved_system}
\end{figure}
The peak in the heat capacity at approximately T=13 indicates the temperature range in which the distribution changes most strongly.
At low temperatures, the system mainly occupies short tours, and the fluctuations in length are small.
With increasing temperature, longer tours become accessible, and the variance increases rapidly.
At high temperatures, the variance approaches a constant value, while $T^2$ continues to increase, causing the heat capacity to decrease again.
The peak therefore marks the temperature range in which the mean tour length changes most strongly with temperature.
It can be seen as a crossover between ordered and disordered tours and should not be seen as a real phase change.
The mean acceptance rate for each temperature is shown in figure \ref{fig:acceptance_annealing_solved}.


\begin{figure}[H]
    \centering
    \includegraphics[width= 0.6\textwidth]{solved_system_acceptance.pdf}
    \caption{Acceptance}
    \label{fig:acceptance_annealing_solved}
\end{figure}
\newpage
\subsubsection{Combined Algorithm}
The mutation in the combined algorithm is done using the Metropolis-Hastings annealing step, as discussed before.
However, the mutation and crossover part of the genetic part breaks detailed balance; therefore, it is not sensible to expect it to sample from the Boltzmann distribution.
Nevertheless, the sampling was compared to the exact solution to observe the difference. 
The evaluation of the combined algorithm was separated into two edge cases, in which it is genetic-dominated and annealing-dominated.
In figure \ref{fig:mixed_solved_system_Nsq_mut}, the annealing-dominated case is seen with the mutation consisting of $N_{\textrm{Nodes}}^2$ proposals.
One can see especially at higher temperatures the observables are sampled from a Boltzmann-like distribution, but at low temperatures there is a mismatch.
This is due to the crossover not being dependent on temperature, while the accepted proposals are. 
Even though the heat capacity is plotted, it should not be interpreted as one, because the genetic part is stochastic and, as said before, breaks detailed balance.
\begin{figure}[H]
    \centering
    \includegraphics[width = 0.7\textwidth]{solved_system_0.5_50_0.01_mixed_TSP.pdf}
    \caption{Comparison between the exact solution (red) and the sampled observables using the mixed algorithm at $N_{\textrm{Nodes}}^2$ annealing proposals per sweep}
    \label{fig:mixed_solved_system_Nsq_mut}
\end{figure}

In figure \ref{fig:mixed_solved_system_4mut}, the $N_{\textrm{Nodes}}^2$ proposals, or 'mutations', are restricted every 5th sweep. 
Now the selection and crossover part dominates for an interval of 5 sweeps, which results in a high mismatch between the exact Boltzmann solution and the generated length means.
At low temperature the mismatch has the same cause as in \ref{fig:mixed_solved_system_Nsq_mut}, while at higher temperatures the crossover and selection show the tendency to restrict the routes with higher lengths.
Even when the acceptance of the annealing algorithm is between 0.7 and 0.8 at T=50, the full $N_{\textrm{Nodes}}^2$ proposals every 5 sweeps don't allow for longer lengths to be generated.
The selection and the crossover, on the other hand, seem to produce shorter lengths on average and are, as discussed before, not temperature dependent, as there is no rejection based on the temperature.
This is expected and shows the stochastic nature of this optimizer fully. 
\begin{figure}[H]
    \centering
    \includegraphics[width = 0.7\textwidth]{solved_system_0.5_50_0.01_mixed_TSP_mut121_int5.pdf}
    \caption{Comparison between the exact solution (red) and the sampled observables using the mixed algorithm at $N_{\textrm{Nodes}}^2$  annealing proposals every 5th sweep}
    \label{fig:mixed_solved_system_4mut}
\end{figure}
\newpage
\subsection{Autocorrelation}
Now this is not really intended to be analyzed in the context of the TSP because it is an optimization problem and the goal is to minimize the observable. 
It is not necessary to treat the observable as something that is desired to be sampled, like the magnetization of a spin grid in the Ising model, and therefore is not needed to estimate its error. 
Yet this is a Monte Carlo course, and I wanted to play around with error estimators and effective sampling and maybe gain some useful knowledge for this data set.
In figure \ref{fig:cities_auto} the \code{ch150} configuration from TSPLIB \cite{tsplib} is shown, which is part of a standardized set of TSP problems.
Most of them have a known optimum, which is very useful for benchmarking.
This set itself has an optimal minimum length of 6528. The distance metric in this case is
\begin{equation}
    L(\tau)
=
\sum_{k=0}^{N-1}
\left\lfloor
\sqrt{
(x_{\tau_k}-x_{\tau_{(k+1)\bmod N}})^2
+
(y_{\tau_k}-y_{\tau_{(k+1)\bmod N}})^2
}
+
\frac{1}{2}
\right\rfloor
\end{equation}  
\begin{figure}[H]
    \centering
    \includegraphics[width = 0.5\textwidth]{auto/cities_used.pdf}
    \caption{\code{ch150} data set with optimum = 6528}
    \label{fig:cities_auto}
\end{figure}
The first step was to estimate the quantities from section \ref{sec:validationMC}, like the mean tour length, variance, heat capacity, and the acceptance rate.
For this, 1500 sweeps were done for thermalization, after which 3000 sweeps were performed at constant temperatures. This was done for both algorithms; however, the combined algorithm was initialized with a population size of 16, for which $N^2$ annealing/mutation steps were performed each sweep.
As the combined algorithm was implemented using parallelization, this only took 2x as long as the annealing algorithm, which was quite surprising.
It is again apparent when looking at figure \ref{fig:temperatures_auto} that the combined algorithm doesn't satisfy detailed balance, but this is expected. 
An interesting feature which was observed is that the best lengths from the population of the combined algorithm (lower left plot: mixed) have a similar peak in the heat capacity as the simulated annealing algorithm. 
In this case the best length is just the lower boundary of the lengths present in the population.
Analyzing the heat capacity further for the simulated annealing one can observe a peak between T=(30 - 45), which is useful information as it tells where the lengths themselves change more rapidly.
Therefore, marking a transition between the ordered and disordered states. 
The acceptance rate doesn't differ for the combined and simulated annealing algorithm, indicating limited influence of the selection/crossover on the acceptance of proposals.
Furthermore, for simulated annealing, the temperature for an uphill acceptance of 0.8 was estimated by proposing changes using the move described in methods and tracking the positive changes.
The temperature was estimated to be T=907 for this system using:


\begin{equation}
    T_0 \approx - \frac{\mathrm{median}(\Delta L_+)}{\mathrm{ln} p_o}
\end{equation}

\begin{figure}[H]
    \centering
    \includegraphics[width = 0.8\textwidth]{auto/temperature_observables.pdf}
    \caption{Temperature observables tracked for the selected system. Mixed: Combined algorithm , standard metropolis: metropolis-hastings}
    \label{fig:temperatures_auto}
\end{figure}
\subsubsection{Metropolis-Hastings Algorithm (Annealing)}

The autocorrelation function of the tour length decays over a relatively short number of stored sweeps, but the low-temperature tail and the blocking analysis do not establish a sharp independent-sample regime. 
The estimated correlation time should therefore be regarded as an observable- and estimator-dependent diagnostic rather than proof that the tours themselves are independent.
\begin{figure}[H]
    \centering
    \includegraphics[width = 0.9\textwidth]{auto/selected_acf_and_blocking_standard_replot.pdf}
    \caption{autocorrelation function and blocking analysis for selected temperatures of annealing algorithm}
    \label{fig:acf_standart}
\end{figure}


\subsubsection{Combined Algorithm}
For the mixed method, selection and crossover change the population distribution. The mixed population mean and mixed best are not ordinary single-chain observables. Their short ACF means that the population statistic changes rapidly between generations, but it does not imply independent Boltzmann samples.
\begin{figure}[H]
    \centering
    \includegraphics[width = \textwidth]{auto/selected_acf_and_blocking_mixed_replot.pdf}
    \caption{autocorrelation function and blocking analysis for selected temperatures of combined algorithm }
    \label{fig:acf_mixed}
\end{figure}

%\subsubsection{Autocorrelation time}
%\begin{figure}[H]
%    \centering
%    \includegraphics[width = 0.8\textwidth]{auto/correlation_and_resampling.pdf}
%    \caption{}
%    \label{fig:acf_time_resample}
%\end{figure}



\subsection{Benchmarking}
\subsubsection{ch150}
With the acquired knowledge from the previous section, the temperature curve was selected to start at the calculated $T_0 \approx 1000$.
The combined algorithm was initialized using 16 population members and a mutation interval of 1. 
The temperature curve was selected as follows:
\begin{equation}
    T(\mathrm{sweep}) = 1000\cdot 0.975^{\mathrm{sweep}}
\end{equation} 
To get a fair comparison between the algorithms, the annealing optimization was performed independently 16 times. 
Out of 16 runs, the annealing algorithm's minimal length was 6528, which corresponds to the optimal minimum of the data set.
The minimal distances out of every run were found and averaged over all runs:
\begin{equation}
    \langle L_{min_R,sweeps} \rangle_R = 6615 \pm 46
\end{equation}
while the average sweep count for reaching the minimal distance for a run is: 
\begin{equation}
    \langle \mathrm{sweeps} \rangle_R = 225 \pm 10
\end{equation}
\begin{figure}[H]
    \centering
    \includegraphics[width = 0.8\textwidth]{benchmarking/ch150/annealing/diagnostics.png}
    \caption{Benchmark diagnostics of annealing algorithm for the dataset \code{ch150}}
    \label{fig:ch150_benchmark_annealing}
\end{figure}

The best tour length reached by the combined algorithm was 6549.
The mean of all minimal distances reached by individual population members over all sweeps is:
\begin{equation}
    \langle L_{min_P,sweeps} \rangle_P = 6549 \pm 0
\end{equation}
while the average sweep count for reaching the minimal distance for a population member is: 
\begin{equation}
    \langle \mathrm{sweeps} \rangle_P = 294 \pm 30
\end{equation}

The best tour was reached by every population member in a span of 100 sweeps, indicating small diversity.
\begin{figure}[H]
    \centering
    \includegraphics[width = 0.8\textwidth]{benchmarking/ch150/mixed/diagnostics.png}
    \caption{Benchmark diagnostics of combined algorithm for the dataset \code{ch150}}
    \label{fig:ch150_benchmark_mixed}
\end{figure}
Looking at Figure \ref{fig:ch150_benchmark_mixed} the impact of the sudden temperature drop is seen in the edge diversity.
At first there are almost no shared edges, which is then reduced, notably with a lag, after the temperature drops.
For low temperatures regimes the edge diversity drops quite rapidly and new beneficial mutations due to the low temperature can only contribute restrictivly.
\begin{figure}[H]
    \centering
    \includegraphics[width = 0.8\textwidth]{benchmarking/ch150/mixed/pop_index_history_detailed.png}
    \caption{Populations family tree in the vicinity of the best tour}
    \label{fig:ch150_benchmark_tree}
\end{figure}
When looking at the family tree vicinity of the best found tour in Figure \ref{fig:ch150_benchmark_tree} one can observe that the best population member was obtained after the crossover of two different population members.
But after two sweeps this best tour was destroyed through two consecutive mutations.
\polyfig{benchmarking/ch150/annealing/final_tour.png}{Best tour found by simulated annealing }{fig:annealing_ch150_besttour}{benchmarking/ch150/mixed/best_tour_overall.png}{Best tour found by combined algorithm }{}
\subsubsection{dsj1000 - temperature cycling}
Additionally, the \code{dsj1000} data set was used to benchmark the optimizers. 
This dataset has 1000 nodes, which makes it computationally more expensive but more interesting to optimize, because the state space is substantially larger. 
It also gives the opportunity to look more deeply into the population behavior. \\ \\
\noindent
First we look at the simulated annealing optimizer.
As both optimizers allow for using custom temperature functions, a temperature cycling approach was performed.
The temperature function consists of an exponential decay followed by two Gaussian distributions. 
This should disorganize the current state space vector to allow it to escape from a local minimum. \\ \\ 

\noindent 
The combined algorithm was initialized with a population of 20 members and a mutation interval of once every 5th sweep.
The mutation itself was done using $N_{\mathrm{Nodes}}^2$ proposals. 
Again,to get a fair comparison between the two algorithms, the annealing algorithm was run consecutively 20 times, adding the run index to the starting seed to get 20 independent trajectories.
Both algorithms were run with 1500 sweeps.
Out of 20 runs the annealing algorithms minimal lenght was 18814218 which it reached at sweep number 946.
The minimal distances out of every run where found and averaged over all runs:
\begin{equation}
    \langle L_{min_R,sweeps} \rangle_R = 19049045 \pm 103388
\end{equation}
while the average sweep count for reaching the minimal distance for a run is: 
\begin{equation}
    \langle \mathrm{sweeps} \rangle_R = 633 \pm 320
\end{equation}
\begin{figure}[H]
    \centering
    \includegraphics[width = 0.8\textwidth]{benchmarking/annealing/diagnostics.png}
    \caption{Benchmark diagnostics of annealing algorithm for the dataset \code{dsj1000}}
    \label{fig:dsj1000_benchmark_annealing}
\end{figure}


The best tour length reached by the combined algorithm was 18955936.
The mean of all minimal distances reached by individual population members over all sweeps is:
\begin{equation}
    \langle L_{min_P,sweeps} \rangle_P = 18967378 \pm 3967
\end{equation}
while the average sweep count for reaching the minimal distance for a population member is: 
\begin{equation}
    \langle \mathrm{sweeps} \rangle_P = 763 \pm 89
\end{equation}
The small standard deviation in the mean of the minimal distances as well as the sweep count at which this length is reached indicates that the population was trapped in a basin consisting of similar local minima.
The diversification did not reach acceptable highs to escape this minima, most likely becuase the mutation step didnt accept enough uphill proposals. 
Figure \ref{fig:dsj1000_benchmark_mixed} supports this claim as in the region the minima was reached the temperature was low. 
The edge diversity was at its low and the amount of shared subsequences larger than 100 was quite high.
The basin itself was most likely reached due to the temperature cycling as it gave the population enough room for diversification.
The minimal distance of all population members was reached 200 sweeps later, directly before the second cycle.
\begin{figure}[H]
    \centering
    \includegraphics[width = 0.8\textwidth]{benchmarking/mixed/diagnostics.png}
    \caption{Benchmark diagnostics of combined algorithm for the dataset \code{dsj1000}}
    \label{fig:dsj1000_benchmark_mixed}
\end{figure}
When looking at the populations family tree in the region of minimal length an interesting observation can be made.
There are two prominent population members, at index 2 and 15 with lenghts 18968182 and 18998150 respectivly, which evade death over longer sweep counts. 
These are most likely part of the best members in the population.
At sweep 945, just before the minimal distance was reached both population members mate generating the population member 5. 
As every 5th sweep there was a mutation step, this member was mutated.
This caused the length of this member to dicrease to the new population minimum. 
It seems that the genetic information of the two top 5 population members was beneficial for the metropolis mutation. 
Before the mutation the crossover of population member 2 ($L_{945}$ = 18968182) and population member 15 ($L_{945}$ = 18998150) created a tour of length 18968527.
\begin{figure}[H]
    \centering
    \includegraphics[width = 0.8\textwidth]{benchmarking/mixed/pop_index_history_detailed.png}
    \caption{Populations family tree in the vicinity of the best tour}
    \label{fig:dsj1000_population_detailed}
\end{figure}
In summary the annealing algorithm found a better minimum than the combined algorithm which is approximately 0.7 \% lower in length.
However the average minimal distance of annealing is higher and has a bigger standart deviation indicating that this was most likely a lucky run. 
Also the amount of sweeps needed to reach a minimal length varies quite strong.
On the other hand the combined algorithm average minimal length spread is lower with an lower average overall.
This indicates that the crossover and selection indeed have a positive influence, but a larger dataset with bigger coverage of the parameters aswell as different seeds are needed to confirm the trend.

\polyfig{benchmarking/annealing/final_tour.png}{Best tour found by simulated annealing using temperature cycling}{fig:annealing_dsj1000_besttour}{benchmarking/mixed/best_tour_overall.png}{Best tour found by combined algorithm using temperature cycling}{}